\begin{figure}[ht]
\centering
\begin{tikzpicture}[x=1cm, y=1cm, font=\small]

% Axes: x = log10(N), y = partial sum value
\draw[->, thick] (-0.2, 0) -- (10.5, 0) node[right] {$N$};
\draw[->, thick] (0, -0.2) -- (0, 5.5) node[above] {$S_N$};

% x-axis ticks: log10 N from 1 to 6
\foreach \x/\lbl in {1/$10$, 2/$10^{2}$, 3/$10^{3}$, 4/$10^{4}$, 5/$10^{5}$, 6/$10^{6}$, 7/$10^{7}$, 8/$10^{8}$, 9/$10^{9}$} {
  \draw (\x, -0.1) -- (\x, 0.1) node[below, yshift=-0.15cm] {\lbl};
}
% y-axis ticks
\foreach \y in {1, 2, 3, 4, 5} {
  \draw (-0.1, \y) -- (0.1, \y) node[left, xshift=-0.15cm] {\y};
}

% Harmonic series: S_N ~ ln N + gamma. Approximate at decades.
% N=10: 2.929; N=100: 5.187; would exceed plot; so scale by 1/3
% Let's plot S_N/4 to fit; relabel y-axis accordingly
% Simpler: plot the geometric-series partial sum sum 0.5^n = 2 - 2*0.5^N -> converges to 2
% and alternating harmonic -> log 2 = 0.693
% and harmonic -> diverges

% Geometric r=0.5 sum: 2(1 - 0.5^N), bound 2
\draw[thick, engblue] plot coordinates {
  (0.3,1.0) (1,1.998) (2,2.0) (3,2.0) (4,2.0) (5,2.0) (6,2.0) (7,2.0) (8,2.0) (9,2.0)
};
\draw[dashed, engblue] (0, 2) -- (9.5, 2) node[right, engblue] {$=2$};
\node[engblue, anchor=west] at (4, 2.25) {geometric $r{=}0.5$};

% Alternating harmonic: -> log 2 = 0.693
\draw[thick, engred] plot coordinates {
  (0.3,0.583) (1,0.646) (2,0.688) (3,0.6927) (4,0.6931) (5,0.6931) (6,0.6931) (9,0.6931)
};
\draw[dashed, engred] (0, 0.693) -- (9.5, 0.693) node[right, engred] {$=\log 2$};
\node[engred, anchor=west] at (5.5, 0.4) {alternating harmonic};

% Harmonic divergence: log N + gamma; scale: plot value/4 so visible. Show through y=5
\draw[thick, engblack] plot coordinates {
  (0.3,0.155) (1,0.733) (2,1.297) (3,1.872) (4,2.448) (5,3.024) (6,3.601) (7,4.177) (8,4.753) (9,5.329)
};
\node[engblack, anchor=west] at (6.5, 4.5) {harmonic};
\node[engblack, anchor=west, font=\scriptsize] at (6.5, 4.1) {(divergent)};

\end{tikzpicture}
\caption[Partial sums of three canonical series]{Partial sums on a
semi-log $N$ axis for three series with shrinking terms. The geometric
series with $r = 0.5$ settles to $2$ after a handful of terms; the
alternating harmonic series settles to $\log 2$ after about a hundred;
the harmonic series grows as $\log N + \gamma$ and never settles. All
three have $a_n \to 0$, yet only two converge. The visual lesson
matches the section 4.6 principle: shrinking terms are necessary but
not sufficient for convergence. The harmonic curve passes $5$ at
$N \approx 10^{9}$.}
\label{fig:vol02:ch04:partial-sums}
\end{figure}
